\section{Results}
\label{sec:results}

All numbers below are computed over the completed corpus: 21,570 of
21,600 planned traces (99.86\%; 30 permanently missing for
Llama-3.3-70B-Instruct only, blocked by an Azure content-safety
filter---see Limitations), 216,319 labeled derived memories, 3 models,
30 scenarios (24 poisoned, 6 clean controls). Every reported point
estimate for H1--H4 carries a 95\% bootstrap CI (10,000 resamples,
stratified by \texttt{poison\_form}, \texttt{scenario\_id} as the
resampling unit, n=24; see \texttt{results/metrics/h\{1,2,3,4\}\_bootstrap\_ci.csv}).

\subsection{H1---Recall vs.\ inflation growth}

Context-exposure (conservative) propagation holds recall at 1.000 by
construction across all three models and every depth---expected, not a
finding on its own. The finding is what it costs: precision degrades
and inflation grows with depth, consistently across all three models
(pooled across 24 scenarios, 95\% CI in brackets; see also
Figure~\ref{fig:precision-vs-depth} and Figure~\ref{fig:blast-radius-growth}):

\begin{table}[h]
\centering
\caption{H1: precision and inflation, context-exposure policy.}
\label{tab:h1}
\begin{tabular}{lcll}
\toprule
Model & Depth & $P_{BR}$ & Inflation \\
\midrule
gpt-4o-mini & 1 & 0.669 [0.642, 0.700] & 1.832 [1.748, 1.910] \\
gpt-4o-mini & 5 & 0.622 [0.593, 0.649] & 2.273 [2.044, 2.516] \\
gpt-4o & 1 & 0.628 [0.609, 0.643] & 1.946 [1.904, 1.995] \\
gpt-4o & 5 & 0.581 [0.552, 0.607] & 2.422 [2.100, 2.783] \\
llama-3.3-70b & 1 & 0.632 [0.620, 0.647] & 1.939 [1.895, 1.975] \\
llama-3.3-70b & 5 & 0.586 [0.552, 0.614] & 2.368 [2.034, 2.773] \\
\bottomrule
\end{tabular}
\end{table}

Structural-only propagation is the mirror image: near-perfect precision
(0.92--1.00) but recall stuck below 1.0 (0.91--0.98), \emph{not}
improving with depth---the \texttt{CO\_RETRIEVED} under-detection
failure mode (a memory merely co-retrieved into a write's context, not
a declared structural parent, whose content nonetheless surfaces the
harmful claim) reproduces at full scale, not just in the original
48-trace pilot that first surfaced it.

\textbf{Every structural-vs-context\_exposure difference is
statistically significant} (Wilcoxon signed-rank, paired by
\texttt{scenario\_id}, n=24, depth 5): $p < 0.001$ for $P_{BR}$,
$R_{BR}$, and inflation\_ratio, across all three models
($p \approx 1.2 \times 10^{-7}$ for $P_{BR}$ and inflation in every
model). The precision/recall/inflation tradeoff is real, not sampling
noise.

\subsection{H2---Attribution-threshold frontier}

Thresholding by cosine similarity between parent/child embeddings
produces a real frontier, not a collapsed scalar (pooled across 24
scenarios, depth 5; see also Figure~\ref{fig:precision-recall-frontier}):

\begin{table}[h]
\centering
\caption{H2: attribution-threshold frontier, depth 5, 95\% CI.}
\label{tab:h2}
\begin{tabular}{llll}
\toprule
Model & Threshold & $P_{BR}$ [95\% CI] & $R_{BR}$ [95\% CI] \\
\midrule
gpt-4o-mini & 0.50 & 0.898 [0.868, 0.930] & 0.990 [0.976, 0.999] \\
gpt-4o-mini & 0.70 & 0.820 [0.781, 0.860] & 0.871 [0.828, 0.916] \\
gpt-4o-mini & 0.85 & 0.633 [0.566, 0.701] & 0.632 [0.558, 0.711] \\
gpt-4o & 0.50 & 0.887 [0.841, 0.928] & 0.998 [0.995, 1.000] \\
gpt-4o & 0.70 & 0.892 [0.844, 0.935] & 0.962 [0.941, 0.980] \\
gpt-4o & 0.85 & 0.710 [0.639, 0.786] & 0.725 [0.645, 0.809] \\
llama-3.3-70b & 0.50 & 0.890 [0.845, 0.929] & 0.997 [0.996, 0.999] \\
llama-3.3-70b & 0.70 & 0.878 [0.832, 0.916] & 0.949 [0.921, 0.975] \\
llama-3.3-70b & 0.85 & 0.640 [0.562, 0.720] & 0.652 [0.577, 0.731] \\
\bottomrule
\end{tabular}
\end{table}

A conservative threshold (0.5) recovers nearly all of context-exposure's
recall while already cutting inflation substantially relative to no
threshold at all. An aggressive threshold (0.85) drops recall sharply
(to 0.63--0.73)---confirming H2's specific concern: raising the
threshold prunes exactly the weakly-attributed \texttt{CO\_RETRIEVED}
edges that, per H1, are sometimes the \emph{only} path to real
contamination.

\textbf{Robustness check: does the frontier depend on the embedding
model?} \texttt{all-MiniLM-L6-v2} was an explicitly documented
placeholder, never validated against another embedding space. We
recomputed H2 with a second, architecturally different model
(\texttt{all-mpnet-base-v2}---MPNet, 768-dim, vs.\ MiniLM, 384-dim),
scoped to gpt-4o-mini/seed=0/all 24 poisoned scenarios (25{,}345
memory nodes re-embedded locally; no new generation, no change to the
frozen main-run embeddings). Thresholds are drawn from \emph{this
run's own} similarity distribution (25th/50th/75th percentile), since
cosine similarity has a different scale in a different embedding
space---not the MiniLM run's 0.5/0.7/0.85:

\begin{table}[h]
\centering
\caption{H2 robustness: second embedding model (all-mpnet-base-v2), depth 5.}
\label{tab:h2-second-embedding}
\begin{tabular}{lccc}
\toprule
Threshold (own-dist.\ quantile) & $P_{BR}$ & $R_{BR}$ & Inflation \\
\midrule
0.19 (p25) & 0.684 & 1.000 & 1.93 \\
0.29 (p50) & 0.757 & 1.000 & 1.64 \\
0.44 (p75) & 0.884 & 0.999 & 1.31 \\
\bottomrule
\end{tabular}
\end{table}

The qualitative frontier survives: precision rises and inflation falls
monotonically as the threshold tightens, exactly the direction H2
predicts, though recall barely erodes here (1.000 $\to$ 0.999) versus
MiniLM's sharp drop at its strictest threshold (to 0.63--0.73)---
suggesting \texttt{all-mpnet-base-v2} separates
\texttt{STRUCTURAL\_PARENT} from \texttt{CO\_RETRIEVED} edges more
cleanly at the high end, not that the frontier is embedding-specific.

\subsection{H3---Surface vs.\ structural traceability (non-directional, as pre-registered)}

Overall laundering rate across the full corpus: \textbf{0.79\%} (662 of
84,177 CARRIES-labeled memories)---small, but real, and this revises
our own pilot's ``zero organic laundering'' finding, which we report
alongside this result rather than silently replacing (per the
deviation-log protocol). Per-model, with 95\% bootstrap CI (24
scenarios): gpt-4o 0.41\% [0.04\%, 1.04\%], gpt-4o-mini 0.46\% [0.00\%,
0.95\%], llama-3.3-70b 1.37\% [0.55\%, 2.44\%]---llama-3.3-70b's CI
does not overlap the other two models', the first bootstrap-confirmed
statistical signal for the cross-model asymmetry this section reports
qualitatively below (see also Figure~\ref{fig:laundering-rate}).
Laundering is highly non-uniform:
\begin{itemize}
\item Concentrated almost entirely in \texttt{continue} and
  \texttt{paraphrase} derivation transforms; \texttt{summarize} and
  \texttt{refine} remain at or near 0\% for all three models,
  consistent with the pilot's original observation for those specific
  transforms.
\item \textbf{A real cross-model asymmetry}: Llama-3.3-70B's
  \texttt{paraphrase} output launders at 4.7--5.9\% by depth 2--5,
  while gpt-4o and gpt-4o-mini's \texttt{paraphrase} output launders at
  0\% throughout. gpt-4o-mini's \texttt{continue} transform launders at
  a measurable, roughly depth-stable 1.0--2.2\%; gpt-4o's
  \texttt{continue} transform launders at a smaller, depth-growing
  0.4--1.7\%.
\end{itemize}

Per H3's non-directional pre-registration, this is reported as
observed: laundering is rare in aggregate, model- and
transform-dependent in a way a single-model pilot could not show, and
does not support a simple ``laundering increases with depth'' story---several
transform/model combinations show no depth trend at all.

\subsection{H4---Containment cost}

\texttt{depth\_aware} (conservative propagation within 2 hops of the
compromised root, structural-only beyond it) reduces over-quarantine
relative to \texttt{flat\_transitive} (unconditional conservative
propagation) at a quantifiable, non-zero recall cost:

% table* (not table): the fourth column made this too wide for a single
% IEEEtran column and its text was overflowing into the adjacent column --
% spanning both columns at the top of the page is the standard IEEE fix.
\begin{table*}[!t]
\centering
\caption{H4: containment cost at depth 5.}
\label{tab:h4}
\begin{tabular}{lccc}
\toprule
Model & Flagged (flat) & Flagged (depth\_aware) & Missed contam.\ (depth\_aware) \\
\midrule
gpt-4o-mini & 10.0 & 7.0 & 0.42 \\
gpt-4o & 10.0 & 7.0 & 0.11 \\
llama-3.3-70b & 10.0 & 7.0 & 0.09 \\
\bottomrule
\end{tabular}
\end{table*}

\texttt{flat\_transitive} misses essentially nothing (missed
contamination $\approx$0.0005 objects/trace, gpt-4o-mini only, zero for
the other two models)---it flags almost everything, so it rarely
misses real contamination, at the cost of the largest quarantine set.
\texttt{depth\_aware} flags \textbf{30\% fewer objects} at depth 5
across all three models, but the recall it gives up is markedly uneven
across models: gpt-4o-mini pays roughly 4--5x the missed-contamination
cost that gpt-4o and llama-3.3-70b do for the identical object-count
reduction---the same real cross-model asymmetry H3 surfaces, showing up
again in a different metric. The \texttt{depth\_aware} vs.\
\texttt{flat\_transitive} gap in missed-contamination count is
statistically significant at every depth $\geq$3 for all three models
(Wilcoxon signed-rank paired by \texttt{scenario\_id}, n=24: e.g.\
gpt-4o-mini depth 5 $p = 0.00013$, gpt-4o depth 5 $p = 0.00044$,
llama-3.3-70b depth 5 $p = 0.00020$; full table in
\texttt{results/metrics/h4\_wilcoxon.csv}).

\textbf{Measurement note, reported rather than hidden:} our original
$C_{regen}$ operationalization (distinct depths touched) does not
differentiate the two policies in this harness's topology, since the
always-continuing structural lineage guarantees every depth is touched
regardless of policy---see Limitations. The object-count numbers above
are the metric that actually demonstrates the tradeoff.

\textbf{Robustness check: does the tradeoff depend on the specific
2-hop window?} \texttt{depth\_aware}'s window was a dated,
explicitly-flagged operationalization (Limitations); a natural
question is whether the reported $\sim$30\% reduction is a real point
on a continuous tradeoff or a hand-picked favorable one. We swept
\texttt{depth\_aware\_window} $\in \{1, 2, 3\}$ at depth 5 (full
24-scenario corpus, all 3 models, all 5 seeds, $n=24$ scenarios),
reusing \texttt{b\_flagged} unmodified---no new code in the frozen
metrics module, only a new caller:

\begin{table}[h]
\centering
\caption{H4 robustness: depth-window sensitivity at depth 5.}
\label{tab:h4-window}
\begin{tabular}{lcc}
\toprule
Window (hops) & Objects flagged (mean) & Missed contam.\ (mean) \\
\midrule
1 & 6.00 & 0.278 \\
2 & 7.00 & 0.206 \\
3 & 8.00 & 0.141 \\
\bottomrule
\end{tabular}
\end{table}

The tradeoff moves smoothly as the window expands: each hop costs
exactly one more flagged object (mechanical, given this harness's
always-continuing structural lineage) and buys back a further
$\sim$25--32\% reduction in missed contamination. This is the stronger
result---not that window=2 is optimal, but that \texttt{depth\_aware}
is one point on a family of depth-dependent policies with a continuous,
predictable tradeoff, not a hand-picked configuration.

\subsection{Clean-control false-positive baseline}

Six genuinely unpoisoned scenarios (no adversarial content anywhere)
provide a false-positive baseline the original 24-scenario design
lacked. Two independent signals (see also
Figure~\ref{fig:clean-control}):

\textbf{Labeler false-positive rate: 0.0000\%} (0 of 43,200 derived
memories labeled CARRIES/REFERENCES against a target that was never
actually present)---perfect specificity on content with no adversarial
signal at all, a meaningful validation of the labeling pipeline
independent of the $\kappa=0.87$--$0.88$ human-agreement result.
(Revised 2026-09-21: a human review of the 6 \texttt{clean\_control}
scenarios found 3 (\texttt{\_03}/\texttt{\_05}/\texttt{\_06}) had
lexical/topical proximity between a benign source fact and the
scenario's harmful semantic target strong enough to risk exactly this
false-positive pattern; those scenarios' benign facts were reworded to
remove the shared vocabulary and the corpus regenerated. The original
rate, 0.0046\% (2/43,200), is reported here as the record of that
correction, not silently overwritten---see
\texttt{configs/scenarios/clean\_control\_review.md}.)

\textbf{Propagation false-positive count} (objects flagged when
$B_{true}=0$ by construction, confirmed for all 86,400 raw rows):
\texttt{structural} flags only the single unavoidable
continuing-lineage node per depth (1, 2, 3, 4, 5---zero variance);
\texttt{context\_exposure}/\texttt{flat\_transitive} flag substantially
more purely from distractor co-retrieval fan-out (10 by depth 5,
consistent across all three models); \texttt{depth\_aware} sits at 7
by depth 5---the same $\sim$30\% reduction seen in H4's
missed-contamination tradeoff, now also confirmed on a corpus with zero
real contamination to miss, not just on the poisoned corpus. This is
independent evidence that \texttt{depth\_aware}'s over-quarantine
reduction is a real structural property of the policy, not an artifact
of the specific poisoned scenarios it was measured on.

\subsection{Real-document ecological-validity slice}
\label{sec:real-document-slice}

The strongest evidence that H1--H4 are not artifacts of a synthetic,
LLM-authored corpus comes from moving the source material out from
under the pipeline: a held-out validation slice from
20 frozen public documents spanning five domains (NIST, CISA, OWASP, FTC,
AWS, Azure, FDA, IRS). We injected one controlled target claim per document
while retaining verbatim benign source facts and distractors, and ran the
frozen generation and labeling pipeline without prompt, labeler, threshold,
or traversal changes. Across 240 traces and 2{,}400 derived memories, the
principal containment tradeoff reproduced: \texttt{depth\_aware} propagation
reduced inflation from 3.57 [2.62, 4.72] under \texttt{flat\_transitive}
propagation to 2.50 [1.82, 3.31] ($n=20$ scenarios as the resampling unit,
mirroring the main corpus's \texttt{scenario\_id}-level bootstrap
discipline), while \texttt{structural}-only propagation achieved
$P_{BR}=0.796$ [0.689, 0.890] and $R_{BR}=0.973$ [0.960, 0.985]---broadly
consistent with the synthetic corpus's precision/recall tradeoff.

Clean-sibling contamination also persisted outside the synthetic corpus:
across the 20 scenarios, the \texttt{child\_2} branch (whose true
structural parents exclude the injected fact) showed strict
semantic-leakage rate 3.83\% [2.00\%, 5.83\%] and broad
(CARRIES+REFERENCES) rate 15.83\% [9.92\%, 22.00\%]. We manually read the
full population of 46 \texttt{child\_2=CARRIES} memories (not a sample)
rather than assuming the mechanism: 40/46 (87\%) explicitly contained the
injected claim's marker or wording despite lacking a structural-parent
edge to it, 5/46 were ambiguous, and 1/46 looked like a labeler
over-trigger---direct textual confirmation that this is genuine
co-retrieval leakage, not a label-rate artifact.

Surface-marker laundering (H3) was substantially higher than in the
synthetic corpus and showed marked scenario-level heterogeneity: the
scenario-level rate ($n=20$, the primary estimand, same resampling
discipline as above) was 19.1\% [6.9\%, 33.1\%], against 0.79\% in the
synthetic corpus. Most individual scenarios laundered 0\% of their CARRIES
memories, while four laundered 30--97\%; the pooled-over-memories figure
(13.5\%, 128/946) understates this heterogeneity by overweighting the
highest-$n$ scenarios. We report the wide CI plainly rather than treating
the point estimate alone as a corpus-invariant quantity: the direction of
the real-vs-synthetic gap is a robust finding, but its exact magnitude is
sensitive to source-text regime and should not be extrapolated to
deployment settings from this slice alone.

This is a more credible claim than ``everything replicated'': the
provenance/containment frontier appears robust across source regimes,
while the absolute surface-laundering rate does not---and the manual
read of all 46 cases shows that difference is measured, not merely
inferred from label rates.

\subsection{Summary across hypotheses}

The consistent thread across H1, H3, and H4 is a genuine cross-model
asymmetry that a single-model study (the original pilot, and most
comparable prior work) could not surface: gpt-4o-mini shows the
steepest H1 inflation growth and the largest H4 missed-contamination
cost under \texttt{depth\_aware}; llama-3.3-70b shows the most
pronounced H3 laundering under \texttt{paraphrase}. No model is
uniformly ``safer'' across all three hypotheses---which model looks
best depends on which failure mode (over-tainting, missed
contamination, or laundering) an operator weights most heavily, itself
a result worth stating plainly rather than collapsing into a single
ranking.

\begin{figure}[t]
\centering
\includegraphics[width=\columnwidth]{../../../results/figures/fig1_precision_vs_depth_by_topk.png}
\caption{Precision vs.\ depth, by \texttt{top\_k}.}
\label{fig:precision-vs-depth}
\end{figure}

\begin{figure}[t]
\centering
\includegraphics[width=\columnwidth]{../../../results/figures/fig2_precision_recall_frontier.png}
\caption{H2 precision/recall frontier under attribution thresholding.}
\label{fig:precision-recall-frontier}
\end{figure}

\begin{figure}[t]
\centering
\includegraphics[width=\columnwidth]{../../../results/figures/fig3_laundering_rate.png}
\caption{H3 laundering rate by model and derivation transform.}
\label{fig:laundering-rate}
\end{figure}

\begin{figure}[t]
\centering
\includegraphics[width=\columnwidth]{../../../results/figures/fig4_blast_radius_growth.png}
\caption{Blast-radius growth with depth.}
\label{fig:blast-radius-growth}
\end{figure}

\begin{figure}[t]
\centering
\includegraphics[width=\columnwidth]{../../../results/figures/fig5_clean_control_false_positive.png}
\caption{Clean-control false-positive baseline.}
\label{fig:clean-control}
\end{figure}
